\subsection{Diagrama de arquitectura}

La arquitectura que se ha desarrollado permite comparar la comunicación
entre procesos mediante sockets TCP y gRPC, cada nodo se ejecuta
como un proceso independiente y dispone de un cliente, un servidor
y un reloj lógico de Lamport compartido entre sus componentes.

Los nodos se ejecutan en el entorno local (\texttt{localhost}),
utilizando los puertos 5000, 5001, 5002 y 5003, los tres primeros
participan en el escenario principal, mientras que el cuarto
se incorpora cuando se realizan pruebas con cuatro nodos.

Esta misma arquitectura se utiliza para evaluar ambos mecanismos
de comunicación, pero las ejecuciones de TCP y gRPC se realizan
por separado, lo cual permite comparar su comportamiento sin
mantener ambos mecanismos activos simultáneamente.

\begin{figure}[H]
    \centering
    \resizebox{\textwidth}{!}{%
        % Arquitectura de comunicación distribuida: TCP y gRPC
% Requiere \usepackage{tikz} y \usetikzlibrary{arrows.meta,positioning,fit}
\shorthandoff{<>}
\begin{tikzpicture}[
    >=Stealth,
    font=\scriptsize,
    nodo/.style={
        draw,
        rounded corners=3pt,
        thick,
        minimum width=4.2cm,
        minimum height=2.0cm,
        align=center,
        fill=white
    },
    enlace/.style={thick, {Stealth}-{Stealth}},
    opcional/.style={thick, dashed, {Stealth}-{Stealth}}
]

\node[nodo] (A) at (0,0) {
    \textbf{Nodo A}\\
    \texttt{localhost:5000}\\[3pt]
    Cliente\\
    Servidor\\
    Reloj de Lamport compartido
};

\node[nodo] (B) at (8,0) {
    \textbf{Nodo B}\\
    \texttt{localhost:5001}\\[3pt]
    Cliente\\
    Servidor\\
    Reloj de Lamport compartido
};

\node[nodo] (C) at (0,-5) {
    \textbf{Nodo C}\\
    \texttt{localhost:5002}\\[3pt]
    Cliente\\
    Servidor\\
    Reloj de Lamport compartido
};

\node[nodo] (D) at (8,-5) {
    \textbf{Nodo D}\\
    \texttt{localhost:5003}\\[3pt]
    Cliente\\
    Servidor\\
    Reloj de Lamport compartido\\
    \emph{Solo escenario de 4 nodos}
};

\draw[enlace] (A.east) --
    node[above, align=center] {TCP / gRPC\\(ejecuciones separadas)}
    (B.west);

\draw[enlace] (A.south) --
    node[right, pos=0.5, xshift=2mm, align=center]
    {TCP / gRPC\\(ejecuciones separadas)}
    (C.north);

\draw[enlace] (B.south west) --
    node[right, pos=0.55, align=center] {TCP / gRPC\\(ejecuciones separadas)}
    (C.north east);

\draw[opcional] (C.east) --
    node[below, align=center] {Escenario de 4 nodos\\TCP / gRPC}
    (D.west);

\begin{scope}[on background layer]
    \node[
        draw,
        dashed,
        rounded corners=5pt,
        inner sep=12pt,
        fit=(A)(B)(C)(D),
        label={[font=\small]above:Entorno local (localhost)}
    ] {};
\end{scope}

\end{tikzpicture}
\shorthandon{<>}
    }
    \caption{Arquitectura de comunicación entre procesos mediante TCP y gRPC. Fuente: Elaboración propia}
    \label{fig:arquitectura-pe-u1}
\end{figure}

En la Figura~\ref{fig:arquitectura-pe-u1}, las conexiones
continuas bidireccionales representan los enlaces principales entre los nodos,
mientras que la conexión discontinua bidireccional identifica el escenario
adicional de cuatro nodos. Esta última no representa
comunicación asíncrona.